% Part: lambda-calculus
% Chapter: church-rosser
% Section: parallel-beta-eta-reduction

\documentclass[../../../include/open-logic-section]{subfiles}

\begin{document}

\olfileid{lam}{cr}{pbe}

\olsection{اختزال~$\beta\eta$ المتوازي}

المطلوب في هذا القسم إثبات خاصية تشيرش--روسر لاختزال~$\beta\eta$ المتوازي؛
فهو النظير المتوازي لاختزال~$\beta\eta$.

\begin{defn}[اختزال~$\beta\eta$ المتوازي، $\beredpar$] \ollabel{defn:beredpar}
  لعلاقة \emph{اختزال~$\beta\eta$ المتوازي} ($\beredpar$) على الحدود
  التعريف الاستقرائي التالي:
  \begin{enumerate}
    \item \ollabel{defn:beredpar1} ${\OLId{latin-ordinary}{x}} \beredpar {\OLId{latin-ordinary}{x}}$.
    \item \ollabel{defn:beredpar2} من~${\OLId{latin-looped}{N}} \beredpar {\OLId{latin-looped}{N}}'$ يلزم
      $\lambd[{\OLId{latin-ordinary}{x}}][{\OLId{latin-looped}{N}}] \beredpar \lambd[{\OLId{latin-ordinary}{x}}][{\OLId{latin-looped}{N}}']$.
    \item \ollabel{defn:beredpar3} من~${\OLId{latin-looped}{P}} \beredpar {\OLId{latin-looped}{P}}'$ و
      ${\OLId{latin-looped}{Q}} \beredpar {\OLId{latin-looped}{Q}}'$ يلزم~$PQ \beredpar {\OLId{latin-looped}{P}}'{\OLId{latin-looped}{Q}}'$.
    \item \ollabel{defn:beredpar4} من~${\OLId{latin-looped}{N}} \beredpar {\OLId{latin-looped}{N}}'$ و
      ${\OLId{latin-looped}{Q}} \beredpar {\OLId{latin-looped}{Q}}'$ يلزم
      $(\lambd[{\OLId{latin-ordinary}{x}}][{\OLId{latin-looped}{N}}]){\OLId{latin-looped}{Q}} \beredpar \Subst{{\OLId{latin-looped}{N}}'}{{\OLId{latin-looped}{Q}}'}{{\OLId{latin-ordinary}{x}}}$.
    \item \ollabel{defn:beredpar5} من~${\OLId{latin-looped}{N}} \beredpar {\OLId{latin-looped}{N}}'$ يلزم
      $\lambd[{\OLId{latin-ordinary}{x}}][Nx] \beredpar {\OLId{latin-looped}{N}}'$، بشرط~${\OLId{latin-ordinary}{x}} \notin \FV{{\OLId{latin-looped}{N}}}$.
  \end{enumerate}
\end{defn}

\begin{thm}\ollabel{thm:refl}
  ${\OLId{latin-looped}{M}} \beredpar {\OLId{latin-looped}{M}}$.
\end{thm}

\begin{proof}
  يترك إثبات ذلك تمرينًا.
\end{proof}

\begin{prob}
  أقم برهان~\olref[lam][cr][pbe]{thm:refl}.
\end{prob}

\begin{defn}[التطوير الكامل بحسب~$\beta\eta$]\ollabel{defn:becd}
  \emph{التطوير الكامل بحسب~$\beta\eta$}، ورمزه~$\becd{{\OLId{latin-looped}{M}}}$، للحد~${\OLId{latin-looped}{M}}$
  هو المحدد بالقواعد الآتية:
  \begin{align}
    \becd{{\OLId{latin-ordinary}{x}}} &= {\OLId{latin-ordinary}{x}} \ollabel{defn:becd1} \\
    \becd{(\lambd[{\OLId{latin-ordinary}{x}}][{\OLId{latin-looped}{N}}])} &= \lambd[{\OLId{latin-ordinary}{x}}][\becd{{\OLId{latin-looped}{N}}}] \ollabel{defn:becd2}
      && \text{إذا لم يكن~${\OLId{latin-looped}{N}}$ من الصورة~${\OLId{latin-looped}{P}} {\OLId{latin-ordinary}{x}}$ مع~${\OLId{latin-ordinary}{x}}\notin\FV{{\OLId{latin-looped}{P}}}$}\\
    \becd{(PQ)} &= \becd{{\OLId{latin-looped}{P}}}\becd{{\OLId{latin-looped}{Q}}} && \text{إذا لم يكن~${\OLId{latin-looped}{P}}$ تجريد~$\lambd$}
    \ollabel{defn:becd3} \\
    \becd{((\lambd[{\OLId{latin-ordinary}{x}}][{\OLId{latin-looped}{N}}]){\OLId{latin-looped}{Q}})} &= \Subst{\becd{{\OLId{latin-looped}{N}}}}{\becd{{\OLId{latin-looped}{Q}}}}{{\OLId{latin-ordinary}{x}}}
                             \ollabel{defn:becd4} \\
    \becd{(\lambd[{\OLId{latin-ordinary}{x}}][Nx])} &= \becd{{\OLId{latin-looped}{N}}} \ollabel{defn:becd5} &
      \text{إذا كان~${\OLId{latin-ordinary}{x}} \notin \FV{{\OLId{latin-looped}{N}}}$}
  \end{align}
\end{defn}

\begin{lem}\ollabel{lem:comp}
  من~${\OLId{latin-looped}{M}} \beredpar {\OLId{latin-looped}{M}}'$ و${\OLId{latin-looped}{R}} \beredpar {\OLId{latin-looped}{R}}'$ يلزم~$\Subst{{\OLId{latin-looped}{M}}}{{\OLId{latin-looped}{R}}}{{\OLId{latin-ordinary}{y}}}
  \beredpar \Subst{{\OLId{latin-looped}{M}}'}{{\OLId{latin-looped}{R}}'}{{\OLId{latin-ordinary}{y}}}$.
\end{lem}

\begin{proof}
  نستقرئ~!!{derivation} العبارة~${\OLId{latin-looped}{M}} \beredpar {\OLId{latin-looped}{M}}'$.
  وعند الحاجة نغير أسماء الممثلين بتحويل~${\OLId{greek-variable}{alpha}}$ قبل النظر في حالات الإحلال،
  حتى يحقق متغير التجريد~${\OLId{latin-ordinary}{x}} \ne {\OLId{latin-ordinary}{y}}$، ولا يقع~${\OLId{latin-ordinary}{x}}$ حرًا في~${\OLId{latin-looped}{R}}$ ولا في~${\OLId{latin-looped}{R}}'$.

  أما الحالات الأربع الأولى فهي بعينها حالات~\olref[pb]{lem:comp}.
  وأما إذا ختم الاشتقاق بالقاعدة~\olref{defn:beredpar5}، فإن~${\OLId{latin-looped}{M}}$
  هو~$\lambd[{\OLId{latin-ordinary}{x}}][Nx]$ و${\OLId{latin-looped}{M}}'$ هو~${\OLId{latin-looped}{N}}'$، لبعض~${\OLId{latin-ordinary}{x}}$ و${\OLId{latin-looped}{N}}$ و${\OLId{latin-looped}{N}}'$،
  مع~${\OLId{latin-ordinary}{x}} \notin \FV{{\OLId{latin-looped}{N}}}$ و${\OLId{latin-looped}{N}} \beredpar {\OLId{latin-looped}{N}}'$.
  فالمطلوب~$\Subst{(\lambd[{\OLId{latin-ordinary}{x}}][Nx])}{{\OLId{latin-looped}{R}}}{{\OLId{latin-ordinary}{y}}} \beredpar \Subst{{\OLId{latin-looped}{N}}'}{{\OLId{latin-looped}{R}}'}{{\OLId{latin-ordinary}{y}}}$،
  أي~$\lambd[{\OLId{latin-ordinary}{x}}][\Subst{{\OLId{latin-looped}{N}}}{{\OLId{latin-looped}{R}}}{{\OLId{latin-ordinary}{y}}} {\OLId{latin-ordinary}{x}}] \beredpar \Subst{{\OLId{latin-looped}{N}}'}{{\OLId{latin-looped}{R}}'}{{\OLId{latin-ordinary}{y}}}$.
  وتحصل هذه النتيجة من~\olref{defn:beredpar}\olref{defn:beredpar5}
  ومن فرض الاستقراء.
\end{proof}

\begin{lem}\ollabel{lem:cont}
  من~${\OLId{latin-looped}{M}} \beredpar {\OLId{latin-looped}{M}}'$ يلزم~${\OLId{latin-looped}{M}}' \beredpar \becd{{\OLId{latin-looped}{M}}}$.
\end{lem}

\begin{proof}
  نستقرئ~!!{derivation} العبارة~${\OLId{latin-looped}{M}} \beredpar {\OLId{latin-looped}{M}}'$.

  تجري الحالات الأربع الأولى على مثال~\olref[pb]{lem:cont}.
  فإن كانت القاعدة الأخيرة~\olref{defn:beredpar5}، كان~${\OLId{latin-looped}{M}}$ هو~$\lambd[{\OLId{latin-ordinary}{x}}][Nx]$
  و${\OLId{latin-looped}{M}}'$ هو~${\OLId{latin-looped}{N}}'$، لبعض~${\OLId{latin-ordinary}{x}}$ و${\OLId{latin-looped}{N}}$ و${\OLId{latin-looped}{N}}'$،
  مع~${\OLId{latin-ordinary}{x}} \notin \FV{{\OLId{latin-looped}{N}}}$ و${\OLId{latin-looped}{N}} \beredpar {\OLId{latin-looped}{N}}'$.
  فنطلب~${\OLId{latin-looped}{N}}' \beredpar \becd{(\lambd[{\OLId{latin-ordinary}{x}}][Nx])}$،
  أو، وهو نفسه،~${\OLId{latin-looped}{N}}' \beredpar \becd{{\OLId{latin-looped}{N}}}$؛ وهذا حاصل من فرض الاستقراء مباشرة.
\end{proof}

\begin{thm}\ollabel{thm:cr}
  العلاقة~$\beredpar$ ذات خاصية تشيرش--روسر.
\end{thm}

\begin{proof}
  النتيجة لازمة مباشرة من~\olref{lem:cont}.
\end{proof}

\end{document}
